%=======================================================================
% preambulo.tex — Configuración general del documento
%=======================================================================

%----------------------- Idioma y codificación -------------------------
\usepackage[utf8]{inputenc}
\usepackage[T1]{fontenc}
% En TeX Live/MiKTeX completos y en Overleaf se usa babel-spanish (spanish.ldf).
% El bloque alternativo permite compilar también en instalaciones mínimas.
\IfFileExists{spanish.ldf}
  {\usepackage[spanish,es-tabla,es-nodecimaldot,es-lcroman]{babel}}
  {\usepackage[provide=*,spanish]{babel}}
\usepackage{csquotes}

%----------------------- Tipografía (Times New Roman) ------------------
% mathptmx = Times New Roman en texto Y en matemáticas (es el paquete
% que normalmente se busca cuando se pide "math Times New Roman").
% Funciona en cualquier distribución sin instalar nada extra.
\usepackage{mathptmx}
\usepackage[scaled=0.90]{helvet}   % sans serif compañera (Helvetica/Arial)
\usepackage{courier}               % monoespaciada

% ALTERNATIVA (recomendada si compilas en Overleaf o tienes TeX Live full):
% newtx da un Times de mejor calidad, sobre todo en matemáticas.
% Comenta las tres líneas de arriba y descomenta estas dos:
% \usepackage{newtxtext}
% \usepackage{newtxmath}

\usepackage{microtype}             % micro-ajustes de justificación

%----------------------- Geometría de página ---------------------------
\usepackage[letterpaper,
            top=4cm, bottom=2.5cm,      % guía de normalización UC: 4 cm superior e izquierdo,
            left=4cm, right=2.5cm,      % 2,5 cm derecho e inferior
            headheight=15pt]{geometry}

\usepackage{setspace}
\onehalfspacing                    % interlineado 1.5 (usa \setstretch{1.15} si prefieres)
\setlength{\parskip}{6pt}
\setlength{\parindent}{0pt}        % párrafos separados por espacio, sin sangría (como el PDF)

%----------------------- Paquetes generales ----------------------------
\usepackage{graphicx}
\usepackage{float}
\usepackage[section]{placeins}   % las figuras no cruzan la sección en que se citan
\usepackage{booktabs}
\usepackage{tabularx}
\usepackage{multirow}
\usepackage{longtable}
\usepackage{enumitem}
\usepackage{xcolor}
\usepackage{amsmath,amssymb,amsfonts}
\usepackage{siunitx}
\sisetup{output-decimal-marker={,}, group-separator={.}}
\usepackage{subcaption}
\usepackage[font=small,labelfont=bf]{caption}
\usepackage{tikz}
\usetikzlibrary{arrows.meta, positioning, shapes.geometric, fit, calc,
                decorations.pathreplacing, backgrounds, matrix, chains}

%----------------------- Colores institucionales -----------------------
\definecolor{pucazul}{HTML}{5B9BD5}
\definecolor{pucgris}{HTML}{404040}
\definecolor{miaamarillo}{HTML}{FFC000}

%----------------------- Encabezados y pies ----------------------------
\usepackage{fancyhdr}

% En fancyhdr v4 conviene definir cada estilo con \fancypagestyle para poder
% alternar entre ellos sin perder la configuración.
\fancypagestyle{principal}{%
  \fancyhf{}%
  \fancyhead[L]{\footnotesize\nouppercase{\leftmark}}%
  \fancyhead[R]{\footnotesize\thepage}%
  \renewcommand{\headrulewidth}{0.4pt}%
  \renewcommand{\footrulewidth}{0pt}%
}

% Estilo de las páginas que abren capítulo y de las páginas preliminares
\fancypagestyle{plain}{%
  \fancyhf{}%
  \fancyfoot[C]{\footnotesize\thepage}%
  \renewcommand{\headrulewidth}{0pt}%
  \renewcommand{\footrulewidth}{0pt}%
}

\renewcommand{\chaptermark}[1]{\markboth{#1}{}}
\pagestyle{principal}

%----------------------- Títulos de capítulo y sección -----------------
% Formato "1 Introducción" en una sola línea, como en el PDF original
% (sin la palabra "Capítulo" y sin salto de línea).
\usepackage{titlesec}
\titleformat{\chapter}[hang]
  {\normalfont\LARGE\bfseries\color{pucgris}}
  {\thechapter}{0.6em}{}
\titlespacing*{\chapter}{0pt}{-10pt}{20pt}

\titleformat{\section}[hang]
  {\normalfont\large\bfseries\color{pucgris}}
  {\thesection}{0.6em}{}
\titlespacing*{\section}{0pt}{14pt}{6pt}

\titleformat{\subsection}[hang]
  {\normalfont\normalsize\bfseries\color{pucgris}}
  {\thesubsection}{0.6em}{}
\titlespacing*{\subsection}{0pt}{12pt}{4pt}

%----------------------- Índice ----------------------------------------
% babel redefine estos nombres al inicio del documento, por eso se
% sobrescriben con \AtBeginDocument.
\AtBeginDocument{%
  \renewcommand{\contentsname}{Tabla de contenido}%
  \renewcommand{\bibname}{Referencias}%
  \renewcommand{\listfigurename}{Índice de figuras}%
  \renewcommand{\listtablename}{Índice de tablas}%
  \renewcommand{\tablename}{Tabla}%
}
\setcounter{tocdepth}{2}     % muestra hasta subsecciones (x.y)
% Entradas de los índices con margen derecho elástico (evita líneas desbordadas en el índice de figuras)
\makeatletter\renewcommand{\@tocrmarg}{2.55em plus 1fil}\makeatother
\setlength{\emergencystretch}{1.5em}   % tolerancia extra de justificación con el ancho de texto de 15 cm
\setcounter{secnumdepth}{3}

%----------------------- Bibliografía (APA vía natbib) -----------------
\usepackage[round,authoryear]{natbib}
\bibliographystyle{apalike-es}   % apalike con "y" entre autores
\setlength{\bibsep}{6pt}
% Uso:  \citep{kahl2021birdnet} -> (Kahl et al., 2021)
%       \citet{kahl2021birdnet} -> Kahl et al. (2021)
%
% ALTERNATIVA APA 7 estricta (requiere biber, disponible en Overleaf):
% \usepackage[style=apa,backend=biber]{biblatex}
% \addbibresource{referencias.bib}
% ...y reemplazar \bibliography{referencias} por \printbibliography

%----------------------- Hipervínculos ---------------------------------
\usepackage[
  bookmarks=true, bookmarksnumbered=true,
  colorlinks=true,
  linkcolor=black, citecolor=black,
  urlcolor=pucazul,
  linktocpage=true, pdfstartview=FitH
]{hyperref}
\usepackage{url}
\urlstyle{same}

%=======================================================================
% MARCAS DE REVISIÓN
% Comentarios del profesor supervisor y pendientes propios.
% Para la entrega final basta con cambiar \revisiontrue por \revisionfalse:
% todas las marcas desaparecen sin tocar el texto.
%=======================================================================
\newif\ifrevision
\revisionfalse         % versión limpia (poner \revisiontrue para ver marcas \pendiente/\profe)

\usepackage[textsize=footnotesize]{todonotes}
\setuptodonotes{inline}

% Comentario del profesor (caja amarilla en el texto)
\newcommand{\profe}[1]{\ifrevision%
  \todo[inline,color=yellow!35,bordercolor=orange]{\textbf{Supervisor:} #1}\fi}

% Pendiente propio (caja azul)
\newcommand{\pendiente}[1]{\ifrevision%
  \todo[inline,color=pucazul!25,bordercolor=pucazul]{\textbf{Pendiente:} #1}\fi}

% Marca ligera en línea para "aquí falta una referencia"
\newcommand{\faltacita}{\ifrevision%
  \textsuperscript{\textcolor{red}{\,[cita]}}\fi}

%----------------------- Figuras TikZ anchas ---------------------------
% Encoge el contenido solo si es más ancho que la línea (no agranda los pequeños).
\newsavebox{\fitwidthbox}
\newcommand{\fitwidth}[1]{%
  \sbox{\fitwidthbox}{#1}%
  \ifdim\wd\fitwidthbox>\linewidth
    \resizebox{\linewidth}{!}{\usebox{\fitwidthbox}}%
  \else
    \usebox{\fitwidthbox}%
  \fi}

%----------------------- Atajos útiles ---------------------------------
\newcommand{\etal}{\textit{et al.}}
\newcommand{\ie}{\textit{i.e.}}
\newcommand{\eg}{\textit{e.g.}}
\DeclareMathOperator*{\argmax}{arg\,max}
\DeclareMathOperator*{\argmin}{arg\,min}

% Evita "Capítulo" en referencias cruzadas mal formateadas
\newcommand{\capref}[1]{Capítulo~\ref{#1}}
\newcommand{\secref}[1]{Sección~\ref{#1}}
\newcommand{\figref}[1]{Figura~\ref{#1}}
\newcommand{\tabref}[1]{Tabla~\ref{#1}}

%----------------------- Ruta de figuras -------------------------------
\graphicspath{{figuras/}}
